% EM structure and teaching revision 20260928
\subsection{开关初值、稳态与自然响应}
理想电感满足 $u_L=L\dot I$，所以在无冲激电压时电流连续；理想电容满足 $I_C=C\dot U_C$，所以在无冲激电流时电压连续。开关刚动作后的初值应从动作前的储能状态继承，不能立即套直流稳态的“电感短路、电容开路”。
串联RLC无源放电，取电容电荷 $Q$ 为状态量，电流方向与 $\dot Q$ 一致，则
\[
 L\ddot Q+R\dot Q+\frac QC=0,
 \qquad s_{1,2}=-\frac R{2L}\pm\sqrt{\frac{R^2}{4L^2}-\frac1{LC}}.
\]
当 $R<2\sqrt{L/C}$ 时为衰减振荡；等号为临界阻尼；较大电阻时两个实负根分别控制快、慢衰减。二阶方程需要两个初值，例如 $Q(0)$ 与 $I(0)$。
若外接恒定电源，先求长期特解，再叠加满足初值的自然响应。检查总能量
$W=Q^2/(2C)+LI^2/2$，无源时应有 $\dot W=-RI^2\le0$。若推导得到无源电阻使振幅持续增长，应首先检查电流参考方向和感应电动势的符号。

